% answers.tex —— 2014 年·湖南卷·文科 · 参考答案与详细解析片段
% 由 tools/generate_projects.py 依来源总账自 2014/2014_数学真题_湖南文_试题.md + 2014/2014_数学真题_湖南文_解析正文.md + 2014/2014_数学真题_湖南文_答案要点.md 生成，勿手改（改源 MD 后重跑）

\begin{answer}
\gksolhead{1}{}

\textbf{题目}：

设命题 $p:\forall x\in\mathbf R,\ x^2+1>0$，则 $\neg p$ 为\mbox{（　　）}
（A）$\exists x_0\in\mathbf R,\ x_0^2+1>0$；（B）$\exists x_0\in\mathbf R,\ x_0^2+1\leqslant0$；（C）$\exists x_0\in\mathbf R,\ x_0^2+1<0$；（D）$\forall x\in\mathbf R,\ x^2+1\leqslant0$。

\Needspace{4\baselineskip}
\textbf{解析}
\Needspace{5\baselineskip}
命题 $p$ 是全称命题。全称命题的否定是特称命题：把量词“$\forall x\in\mathbf R$”改为“$\exists x_0\in\mathbf R$”，并把结论 $x_0^2+1>0$ 改为它的否定 $x_0^2+1\leqslant0$。于是
\nopagebreak[4]
\[\neg p:\ \exists x_0\in\mathbf R,\ x_0^2+1\leqslant0 .\]
对照选项，应选 B。

\gkanswer{答案：B。全称命题的否定是特称命题，量词改为“存在”，结论 $x^2+1>0$ 改为其否定 $x_0^2+1\leqslant0$。}

\end{answer}
\begin{answer}
\gksolhead{2}{}

\textbf{题目}：

已知集合 $A=\{x\mid x>2\}$，$B=\{x\mid1<x<3\}$，则 $A\cap B=$\mbox{（　　）}
（A）$\{x\mid x>2\}$；（B）$\{x\mid x>1\}$；（C）$\{x\mid2<x<3\}$；（D）$\{x\mid1<x<3\}$。

\Needspace{4\baselineskip}
\textbf{解析}
\Needspace{5\baselineskip}
由 $x>2$ 得 $A=(2,+\infty)$；由 $1<x<3$ 得 $B=(1,3)$。在数轴上取交集，公共部分是 $2<x<3$，即
\nopagebreak[4]
\[A\cap B=(2,3)=\{x\mid2<x<3\}.\]
选 C。

\gkanswer{答案：C。$A=(2,+\infty)$，$B=(1,3)$，在数轴上取交得 $(2,3)$，即 $\{x\mid2<x<3\}$。}

\end{answer}
\begin{answer}
\gksolhead{3}{}

\textbf{题目}：

对一个容量为 $N$ 的总体抽取容量为 $n$ 的样本，当选取简单随机抽样、系统抽样和分层抽样三种不同方法抽取样本时，总体中每个个体被抽中的概率分别为 $p_1$，$p_2$，$p_3$，则\mbox{（　　）}
（A）$p_1=p_2<p_3$；（B）$p_2=p_3<p_1$；（C）$p_1=p_3<p_2$；（D）$p_1=p_2=p_3$。

\Needspace{4\baselineskip}
\textbf{解析}
\Needspace{5\baselineskip}
简单随机抽样、系统抽样与分层抽样都是等可能抽样：在容量为 $N$ 的总体中抽取容量为 $n$ 的样本时，每个个体被抽中的机会相同，都等于 $\dfrac nN$。因此
\nopagebreak[4]
\[p_1=p_2=p_3=\frac nN .\]
选 D。

\gkanswer{答案：D。简单随机抽样、系统抽样、分层抽样都是等可能抽样，每个个体被抽中的概率都等于 $\dfrac nN$。}

\end{answer}
\begin{answer}
\gksolhead{4}{}

\textbf{题目}：

下列函数中，既是偶函数又在区间 $(-\infty,0)$ 上单调递增的是\mbox{（　　）}
（A）$f(x)=\dfrac1{x^2}$；（B）$f(x)=x^2+1$；（C）$f(x)=x^3$；（D）$f(x)=2^{-x}$。

\Needspace{4\baselineskip}
\textbf{解析}
逐项判断。
\Needspace{5\baselineskip}
（A）$f(x)=\dfrac1{x^2}$，定义域关于原点对称，且 $f(-x)=\dfrac1{(-x)^2}=f(x)$，是偶函数；当 $x<0$ 时
\nopagebreak[4]
\[f'(x)=-\frac2{x^3}>0,\]
故在 $(-\infty,0)$ 上单调递增，符合题意。
（B）$f(x)=x^2+1$ 是偶函数，但 $x<0$ 时 $f'(x)=2x<0$，在 $(-\infty,0)$ 上单调递减，不符合。
（C）$f(x)=x^3$ 是奇函数，不符合。
（D）$f(x)=2^{-x}$ 满足 $f(-x)=2^{x}\neq f(x)$，不是偶函数，不符合。
选 A。

\gkanswer{答案：A。$\dfrac1{x^2}$ 为偶函数，且 $x<0$ 时导数为 $-\dfrac2{x^3}>0$，在 $(-\infty,0)$ 上递增；$x^2+1$ 在该区间递减，$x^3$ 为奇函数，$2^{-x}$ 不是偶函数。}

\end{answer}
\begin{answer}
\gksolhead{5}{}

\textbf{题目}：

在区间 $[-2,3]$ 上随机选取一个数 $X$，则 $X\leqslant1$ 的概率为\mbox{（　　）}
（A）$\dfrac45$；（B）$\dfrac35$；（C）$\dfrac25$；（D）$\dfrac15$。

\Needspace{4\baselineskip}
\textbf{解析}
\Needspace{5\baselineskip}
在区间 $[-2,3]$ 上随机取数，是长度型几何概型。区间总长度为
\nopagebreak[4]
\[3-(-2)=5,\]
\Needspace{5\baselineskip}
满足 $X\leqslant1$ 的部分是 $[-2,1]$，长度为
\nopagebreak[4]
\[1-(-2)=3.\]
\Needspace{5\baselineskip}
故
\nopagebreak[4]
\[P(X\leqslant1)=\frac35 .\]
选 B。

\gkanswer{答案：B。区间长度为 $5$，满足 $X\leqslant1$ 的部分为 $[-2,1]$，长度为 $3$，故概率为 $\dfrac35$。}

\end{answer}
\begin{answer}
\gksolhead{6}{}

\textbf{题目}：

若圆 $C_1:x^2+y^2=1$ 与圆 $C_2:x^2+y^2-6x-8y+m=0$ 外切，则 $m=$\mbox{（　　）}
（A）$21$；（B）$19$；（C）$9$；（D）$-11$。

\Needspace{4\baselineskip}
\textbf{解析}
圆 $C_1$ 的圆心为 $(0,0)$，半径 $r_1=1$。
\Needspace{5\baselineskip}
把 $C_2$ 的方程配方：
\nopagebreak[4]
\[x^2-6x+y^2-8y+m=0\Longrightarrow(x-3)^2+(y-4)^2=25-m,\]
故圆心为 $(3,4)$，半径 $r_2=\sqrt{25-m}$（需 $m<25$）。
\Needspace{5\baselineskip}
两圆外切时，圆心距等于半径之和：
\nopagebreak[4]
\[\sqrt{(3-0)^2+(4-0)^2}=1+\sqrt{25-m}\Longrightarrow5=1+\sqrt{25-m}.\]
\Needspace{5\baselineskip}
于是 $\sqrt{25-m}=4$，$25-m=16$，得
\nopagebreak[4]
\[m=9 .\]
选 C。

\gkanswer{答案：C。$C_2$ 化为 $(x-3)^2+(y-4)^2=25-m$，圆心距 $5=1+\sqrt{25-m}$，得 $\sqrt{25-m}=4$，故 $m=9$。}

\end{answer}
\begin{answer}
\gksolhead{7}{}

\textbf{题目}：

执行如图所示的程序框图，如果输入的 $t\in[-2,2]$，则输出的 $S$ 属于\mbox{（　　）}
\nopagebreak[4]
\begin{center}% ============================================================
% f2014_hunan_arts_q07_flowchart.tex
% 2014 湖南卷（文科）第 7 题：程序框图
% 来源 PDF：普通高考/2014/2014湖南文.pdf
% 原页：第 1 页（A3 横排）左栏下方，印面 887，第 7 题题干下方
% 结构：开始 → 输入 t → 判断 t<0；
%       是：t=2t^2+1，回到判断框；否：S=t-3 → 输出 S → 结束。
% ============================================================
\begin{tikzpicture}[>=Stealth, line width=0.7pt, font=\small,
  x=1cm, y=1cm, line cap=round, line join=round]
  % 开始
  \draw[rounded corners=6pt] (-0.80,0.35) rectangle (0.80,-0.35);
  \node at (0,0) {开始};
  \draw[->] (0,-0.35) -- (0,-0.62);
  % 输入 t
  \draw (-1.15,-0.62) -- (1.15,-0.62) -- (1.00,-1.28) -- (-1.30,-1.28) -- cycle;
  \node at (0,-0.95) {输入 $t$};
  \draw[->] (0,-1.28) -- (0,-1.58);
  % 判断框
  \draw (0,-1.58) -- (1.95,-2.28) -- (0,-2.98) -- (-1.95,-2.28) -- cycle;
  \node at (0,-2.28) {$t<0$?};
  % 是：向左再向下进入循环体
  \draw[->] (-1.95,-2.28) -- (-2.85,-2.28) -- (-2.85,-3.25);
  \node at (-2.35,-2.12) {是};
  \draw (-3.80,-3.25) rectangle (-1.90,-3.85);
  \node at (-2.85,-3.55) {$t=2t^2+1$};
  % 回边：循环体左侧向上，水平进入判断框左上边
  \draw[->] (-3.80,-3.55) -- (-4.45,-3.55) -- (-4.45,-1.80) -- (-0.61,-1.80);
  % 否：向右再向下
  \draw[->] (1.95,-2.28) -- (3.00,-2.28) -- (3.00,-3.12);
  \node at (2.42,-2.12) {否};
  \draw (2.05,-3.12) rectangle (3.95,-3.72);
  \node at (3.00,-3.42) {$S=t-3$};
  % 输出 S
  \draw[->] (3.00,-3.72) -- (3.00,-4.25) -- (0,-4.25) -- (0,-4.45);
  \draw (-1.35,-4.45) -- (1.35,-4.45) -- (1.05,-5.15) -- (-1.65,-5.15) -- cycle;
  \node at (-0.30,-4.80) {输出 $S$};
  \draw[->] (0,-5.15) -- (0,-5.45);
  % 结束
  \draw[rounded corners=6pt] (-0.80,-5.45) rectangle (0.80,-6.15);
  \node at (0,-5.80) {结束};
\end{tikzpicture}
\end{center}
\nopagebreak[4]
（A）$[-6,-2]$；（B）$[-5,-1]$；（C）$[-4,5]$；（D）$[-3,6]$。

\Needspace{4\baselineskip}
\textbf{解析}
按框图分两种情况。
\Needspace{5\baselineskip}
若 $t<0$，执行 $t\leftarrow2t^2+1$。因为 $2t^2+1\geqslant1>0$，返回判断后条件不再成立，随即执行 $S=t-3$。此时
\nopagebreak[4]
\[S=2t^2+1-3=2t^2-2 .\]
由 $t\in[-2,0)$ 得 $t^2\in(0,4]$，故 $S\in(-2,6]$。
\Needspace{5\baselineskip}
若 $t\geqslant0$，不进入循环体，直接得
\nopagebreak[4]
\[S=t-3\in[-3,-1].\]
\Needspace{5\baselineskip}
两种情况合并：
\nopagebreak[4]
\[S\in[-3,-1]\cup(-2,6]=[-3,6].\]
选 D。

\gkanswer{答案：D。$t<0$ 时先化为 $t=2t^2+1\in[1,9]$，再得 $S=t-3\in[-2,6]$；$t\geqslant0$ 时 $S=t-3\in[-3,-1]$，合并得 $[-3,6]$。}

\end{answer}
\begin{answer}
\gksolhead{8}{}

\textbf{题目}：

一块石材表示的几何体的三视图如图所示，将该石材切削、打磨，加工成球，则能得到的最大球的半径等于\mbox{（　　）}
\nopagebreak[4]
\begin{center}% ============================================================
% f2014_hunan_arts_q08_views.tex
% 2014 湖南卷（文科）第 8 题：石材几何体的三视图
% 来源 PDF：普通高考/2014/2014湖南文.pdf
% 原页：第 1 页（A3 横排）第 2 栏顶部，印面 887，第 8 题题干上方
% 结构：正视图为直角边长 6、8 的直角三角形；侧视图为 12×8 的矩形；
%       俯视图为 6×12 的矩形。比例 0.4 cm/单位。
% ============================================================
\begin{tikzpicture}[>=Stealth, line width=0.6pt, font=\small, scale=0.72,
  x=1cm, y=1cm, line cap=round, line join=round]
  % 正视图：直角边 6（横）、8（纵）
  \draw (0,0) -- (2.4,0) -- (0,3.2) -- cycle;
  \draw[<->] (0,3.95) -- (2.4,3.95);
  \node at (1.2,4.22) {$6$};
  \draw[<->] (-0.95,0) -- (-0.95,3.2);
  \node at (-1.18,1.6) {$8$};
  \node at (1.2,-0.52) {正视图};
  % 侧视图：12×8
  \draw (4.4,0) rectangle (9.2,3.2);
  \node at (6.8,-0.52) {侧视图};
  % 俯视图：6×12
  \draw (10.4,-1.6) rectangle (12.8,3.2);
  \draw[<->] (13.10,-1.6) -- (13.10,3.2);
  \node at (13.52,0.8) {$12$};
  \node at (11.6,-2.12) {俯视图};
\end{tikzpicture}
\end{center}
\nopagebreak[4]
（A）$1$；（B）$2$；（C）$3$；（D）$4$。

\Needspace{4\baselineskip}
\textbf{解析}
由三视图还原几何体：正视图是直角边长分别为 $6$、$8$ 的直角三角形，俯视图是 $6\times12$ 的矩形，侧视图是 $12\times8$ 的矩形。因此该石材是一个直三棱柱，其横截面为两直角边 $6$、$8$、斜边 $10$ 的直角三角形，柱长为 $12$。
\Needspace{5\baselineskip}
球要整个放入柱体，其大圆必须落在三角形截面内，半径不超过该三角形的内切圆半径
\nopagebreak[4]
\[r=\frac{6+8-10}2=2 .\]
又球的直径为 $2r=4<12$，球可置于柱体中部而不超出两端，长度方向不构成限制。故最大球的半径为 $2$。
选 B。

\gkanswer{答案：B。三视图还原为横截面是直角边 $6$、$8$ 的直角三角形、长为 $12$ 的直三棱柱，最大球半径为其内切圆半径 $\dfrac{6+8-10}2=2$。}

\end{answer}
\begin{answer}
\gksolhead{9}{}

\textbf{题目}：

若 $0<x_1<x_2<1$，则\mbox{（　　）}
（A）$\mathrm e^{x_2}-\mathrm e^{x_1}>\ln x_2-\ln x_1$；（B）$\mathrm e^{x_2}-\mathrm e^{x_1}<\ln x_2-\ln x_1$；（C）$x_2\mathrm e^{x_1}>x_1\mathrm e^{x_2}$；（D）$x_2\mathrm e^{x_1}<x_1\mathrm e^{x_2}$。

\Needspace{4\baselineskip}
\textbf{解析}
\Needspace{5\baselineskip}
先看（C）（D）。两式都等价于比较 $\dfrac{\mathrm e^x}{x}$ 的大小：由 $x_1x_2>0$，
\nopagebreak[4]
\[x_2\mathrm e^{x_1}>x_1\mathrm e^{x_2}\iff\frac{\mathrm e^{x_1}}{x_1}>\frac{\mathrm e^{x_2}}{x_2}.\]
\Needspace{5\baselineskip}
令 $g(x)=\dfrac{\mathrm e^x}x$，则
\nopagebreak[4]
\[g'(x)=\frac{\mathrm e^x(x-1)}{x^2}.\]
\Needspace{5\baselineskip}
当 $0<x<1$ 时 $g'(x)<0$，$g$ 在 $(0,1)$ 上单调递减。由 $0<x_1<x_2<1$ 得
\nopagebreak[4]
\[g(x_1)>g(x_2),\]
\Needspace{5\baselineskip}
即 $\dfrac{\mathrm e^{x_1}}{x_1}>\dfrac{\mathrm e^{x_2}}{x_2}$，故
\nopagebreak[4]
\[x_2\mathrm e^{x_1}>x_1\mathrm e^{x_2},\]
（C）成立，（D）不成立。
\Needspace{5\baselineskip}
再排除（A）（B）。固定 $x_2=\dfrac12$，令 $x_1\to0^+$，则
\nopagebreak[4]
\[\ln x_2-\ln x_1\to+\infty,\]
\Needspace{5\baselineskip}
而
\nopagebreak[4]
\[\mathrm e^{x_2}-\mathrm e^{x_1}<\mathrm e^{1/2}<\mathrm e<3,\]
故（A）与（B）都不可能恒成立。
选 C。

\gkanswer{答案：C。令 $g(x)=\dfrac{\mathrm e^x}x$，$g'(x)=\dfrac{\mathrm e^x(x-1)}{x^2}<0$ 于 $(0,1)$，故 $g(x_1)>g(x_2)$，即 $x_2\mathrm e^{x_1}>x_1\mathrm e^{x_2}$。}

\end{answer}
\begin{answer}
\gksolhead{10}{}

\textbf{题目}：

在平面直角坐标系中，$O$ 为原点，$A(-1,0)$，$B(0,\sqrt3)$，$C(3,0)$，动点 $D$ 满足 $|\overrightarrow{CD}|=1$，则 $|\overrightarrow{OA}+\overrightarrow{OB}+\overrightarrow{OD}|$ 的取值范围是\mbox{（　　）}
（A）$[4,6]$；（B）$[\sqrt{19}-1,\sqrt{19}+1]$；（C）$[2\sqrt3,2\sqrt7]$；（D）$[\sqrt7-1,\sqrt7+1]$。

\Needspace{4\baselineskip}
\textbf{解析}
\Needspace{5\baselineskip}
设 $D(x,y)$，由 $|\overrightarrow{CD}|=1$ 得
\nopagebreak[4]
\[(x-3)^2+y^2=1,\]
即 $D$ 在以 $(3,0)$ 为圆心、$1$ 为半径的圆上。
\Needspace{5\baselineskip}
又
\nopagebreak[4]
\[\overrightarrow{OA}+\overrightarrow{OB}=(-1,0)+(0,\sqrt3)=(-1,\sqrt3),\]
\Needspace{5\baselineskip}
故
\nopagebreak[4]
\[\overrightarrow{OA}+\overrightarrow{OB}+\overrightarrow{OD}=(x-1,\ y+\sqrt3).\]
\Needspace{5\baselineskip}
令 $u=x-3$，$v=y$，则 $u^2+v^2=1$，且
\nopagebreak[4]
\[(x-1,\ y+\sqrt3)=(2+u,\ \sqrt3+v).\]
\Needspace{5\baselineskip}
所以所求向量的终点在以 $M(2,\sqrt3)$ 为圆心、$1$ 为半径的圆上，其模的取值范围为
\nopagebreak[4]
\[\bigl||OM|-1,\ |OM|+1\bigr],\qquad |OM|=\sqrt{2^2+(\sqrt3)^2}=\sqrt7 .\]
\Needspace{5\baselineskip}
故取值范围是
\nopagebreak[4]
\[[\sqrt7-1,\ \sqrt7+1].\]
选 D。

\gkanswer{答案：D。$\overrightarrow{OA}+\overrightarrow{OB}=(-1,\sqrt3)$，故所求向量的终点在以 $(2,\sqrt3)$ 为圆心、$1$ 为半径的圆上，模的取值范围为 $[\sqrt7-1,\sqrt7+1]$。}

\end{answer}
\begin{answer}
\gksolhead{11}{}

\textbf{题目}：

复数 $\dfrac{3+\mathrm i}{\mathrm i^2}$（$\mathrm i$ 为虚数单位）的实部等于\gkblank{}。

\Needspace{4\baselineskip}
\textbf{解析}
\Needspace{5\baselineskip}
由 $\mathrm i^2=-1$，
\nopagebreak[4]
\[\frac{3+\mathrm i}{\mathrm i^2}=\frac{3+\mathrm i}{-1}=-3-\mathrm i .\]
其实部为 $-3$。

\gkanswer{答案：$-3$。$\mathrm i^2=-1$，故 $\dfrac{3+\mathrm i}{\mathrm i^2}=-3-\mathrm i$，实部为 $-3$。}

\end{answer}
\begin{answer}
\gksolhead{12}{}

\textbf{题目}：

在平面直角坐标系中，曲线 $C:\begin{cases}x=2+\dfrac{\sqrt2}2t,\\y=1+\dfrac{\sqrt2}2t,\end{cases}$（$t$ 为参数）的普通方程为\gkblank{}。

\Needspace{4\baselineskip}
\textbf{解析}
\Needspace{5\baselineskip}
由参数方程
\nopagebreak[4]
\[x=2+\frac{\sqrt2}2t,\qquad y=1+\frac{\sqrt2}2t\]
\Needspace{5\baselineskip}
两式相减，消去 $t$：
\nopagebreak[4]
\[x-y=(2-1)+\left(\frac{\sqrt2}2-\frac{\sqrt2}2\right)t=1 .\]
\Needspace{5\baselineskip}
故普通方程为
\nopagebreak[4]
\[x-y-1=0 .\]

\gkanswer{答案：$x-y-1=0$。两式相减消去 $t$ 得 $x-y=1$。}

\end{answer}
\begin{answer}
\gksolhead{13}{}

\textbf{题目}：

若变量 $x$，$y$ 满足约束条件 $\begin{cases}y\leqslant x,\\x+y\leqslant4,\\y\geqslant1,\end{cases}$ 则 $z=2x+y$ 的最大值为\gkblank{}。

\Needspace{4\baselineskip}
\textbf{解析}
\Needspace{5\baselineskip}
约束条件
\nopagebreak[4]
\[y\leqslant x,\qquad x+y\leqslant4,\qquad y\geqslant1\]
\Needspace{5\baselineskip}
所围成的可行域是三角形，其三个顶点为三条边界两两的交点：
\nopagebreak[4]
\[y=1\text{ \text{与} }y=x\Longrightarrow(1,1),\]
\nopagebreak[4]
\[y=1\text{ \text{与} }x+y=4\Longrightarrow(3,1),\]
\nopagebreak[4]
\[y=x\text{ \text{与} }x+y=4\Longrightarrow(2,2).\]
\Needspace{5\baselineskip}
分别代入 $z=2x+y$：
\nopagebreak[4]
\[z(1,1)=3,\qquad z(3,1)=7,\qquad z(2,2)=6 .\]
线性目标函数在可行域的顶点处取得最值，故最大值为 $7$。

\gkanswer{答案：$7$。可行域的三个顶点为 $(1,1)$，$(3,1)$，$(2,2)$，代入 $z=2x+y$ 分别得 $3$，$7$，$6$，最大值为 $7$。}

\end{answer}
\begin{answer}
\gksolhead{14}{}

\textbf{题目}：

平面上一机器人在行进中始终保持与点 $F(1,0)$ 的距离和到直线 $x=-1$ 的距离相等。若机器人接触不到过点 $P(-1,0)$ 且斜率为 $k$ 的直线，则 $k$ 的取值范围是\gkblank{}。

\Needspace{4\baselineskip}
\textbf{解析}
\Needspace{5\baselineskip}
到定点 $F(1,0)$ 的距离等于到定直线 $x=-1$ 的距离，故机器人的轨迹是以 $F(1,0)$ 为焦点、$x=-1$ 为准线的抛物线
\nopagebreak[4]
\[y^2=4x .\]
\Needspace{5\baselineskip}
过 $P(-1,0)$ 且斜率为 $k$ 的直线为
\nopagebreak[4]
\[l:\ y=k(x+1).\]
“接触不到”即直线 $l$ 与抛物线没有公共点。
若 $k=0$，则 $l:y=0$ 与抛物线交于顶点 $(0,0)$，能接触到，不合题意。
\Needspace{5\baselineskip}
若 $k\neq0$，代入 $y^2=4x$：
\nopagebreak[4]
\[k^2(x+1)^2=4x\Longrightarrow k^2x^2+(2k^2-4)x+k^2=0 .\]
\Needspace{5\baselineskip}
无公共点等价于判别式小于零：
\nopagebreak[4]
\[\Delta=(2k^2-4)^2-4k^4<0 .\]
\Needspace{5\baselineskip}
化简：
\nopagebreak[4]
\[\begin{aligned}
\Delta&=4(k^2-2)^2-4k^4\\
&=4\left(k^4-4k^2+4-k^4\right)\\
&=16-16k^2<0,
\end{aligned}\]
\Needspace{5\baselineskip}
故 $k^2>1$，即
\nopagebreak[4]
\[k\in(-\infty,-1)\cup(1,+\infty).\]

\gkanswer{答案：$(-\infty,-1)\cup(1,+\infty)$。机器人轨迹为抛物线 $y^2=4x$；直线 $y=k(x+1)$ 与它无公共点，判别式小于零给出 $k^2>1$。}

\end{answer}
\begin{answer}
\gksolhead{15}{}

\textbf{题目}：

若 $f(x)=\ln(\mathrm e^{3x}+1)+ax$ 是偶函数，则 $a=$\gkblank{}。

\Needspace{4\baselineskip}
\textbf{解析}
\Needspace{5\baselineskip}
$f(x)$ 为偶函数，故对任意 $x$ 有 $f(-x)=f(x)$：
\nopagebreak[4]
\[\ln\left(\mathrm e^{-3x}+1\right)-ax=\ln\left(\mathrm e^{3x}+1\right)+ax .\]
\Needspace{5\baselineskip}
而
\nopagebreak[4]
\[\ln\left(\mathrm e^{-3x}+1\right)=\ln\frac{1+\mathrm e^{3x}}{\mathrm e^{3x}}=\ln\left(\mathrm e^{3x}+1\right)-3x .\]
\Needspace{5\baselineskip}
代入得
\nopagebreak[4]
\[\ln\left(\mathrm e^{3x}+1\right)-3x-ax=\ln\left(\mathrm e^{3x}+1\right)+ax,\]
\Needspace{5\baselineskip}
即
\nopagebreak[4]
\[-3x-ax=ax\Longrightarrow-3-a=a\Longrightarrow a=-\frac32 .\]

\gkanswer{答案：$-\dfrac32$。由 $f(-x)=f(x)$ 得 $\ln(\mathrm e^{3x}+1)-3x-ax=\ln(\mathrm e^{3x}+1)+ax$，故 $a=-\dfrac32$。}

\end{answer}
\begin{answer}
\gksolhead{16}{}

\textbf{题目}：

已知数列 $\{a_n\}$ 的前 $n$ 项和 $S_n=\dfrac{n^2+n}2$，$n\in\mathbf N^*$。
（1） 求数列 $\{a_n\}$ 的通项公式；
（2） 设 $b_n=2^{a_n}+(-1)^na_n$，求数列 $\{b_n\}$ 的前 $2n$ 项和。

\Needspace{4\baselineskip}
\textbf{解析}
\Needspace{5\baselineskip}
\textbf{（1）} 当 $n=1$ 时，
\nopagebreak[4]
\[a_1=S_1=\frac{1^2+1}2=1 .\]
\Needspace{5\baselineskip}
当 $n\geqslant2$ 时，
\nopagebreak[4]
\[\begin{aligned}
a_n&=S_n-S_{n-1}\\
&=\frac{n^2+n}2-\frac{(n-1)^2+(n-1)}2\\
&=\frac{n^2+n-\left(n^2-2n+1\right)-(n-1)}2\\
&=\frac{2n}2=n .
\end{aligned}\]
\Needspace{5\baselineskip}
$n=1$ 时 $a_1=1$ 也符合 $a_n=n$，故
\nopagebreak[4]
\[a_n=n\quad(n\in\mathbf N^*).\]
\Needspace{5\baselineskip}
\textbf{（2）} 由（1），
\nopagebreak[4]
\[b_n=2^{a_n}+(-1)^na_n=2^n+(-1)^nn .\]
\Needspace{5\baselineskip}
设 $\{b_n\}$ 的前 $2n$ 项和为 $T_{2n}$，则
\nopagebreak[4]
\[T_{2n}=\sum_{k=1}^{2n}2^k+\sum_{k=1}^{2n}(-1)^kk .\]
\Needspace{5\baselineskip}
等比部分
\nopagebreak[4]
\[\sum_{k=1}^{2n}2^k=\frac{2\left(2^{2n}-1\right)}{2-1}=2^{2n+1}-2 .\]
\Needspace{5\baselineskip}
交错部分两两配对：
\nopagebreak[4]
\[\begin{aligned}
\sum_{k=1}^{2n}(-1)^kk&=(-1+2)+(-3+4)+\cdots+(-(2n-1)+2n)\\
&=\underbrace{1+1+\cdots+1}_{n\text{ \text{个}}}=n .
\end{aligned}\]
\Needspace{5\baselineskip}
故
\nopagebreak[4]
\[T_{2n}=2^{2n+1}-2+n=2^{2n+1}+n-2 .\]

\gkanswer{答案（1）：$a_n=n$。$a_1=S_1=1$，$n\geqslant2$ 时 $a_n=S_n-S_{n-1}=n$，对 $n=1$ 也成立。}

\gkanswer{答案（2）：$2^{2n+1}+n-2$。$b_n=2^n+(-1)^nn$，等比部分求和为 $2^{2n+1}-2$，交错部分求和为 $n$。}

\end{answer}
\begin{answer}
\gksolhead{17}{}

\textbf{题目}：

某企业有甲、乙两个研发小组，为了比较他们的研发水平，现随机抽取这两个小组往年研发新产品的结果如下：$(a,b)$，$(a,\bar b)$，$(a,b)$，$(\bar a,b)$，$(\bar a,\bar b)$，$(a,b)$，$(a,b)$，$(a,\bar b)$，$(\bar a,b)$，$(a,\bar b)$，$(\bar a,\bar b)$，$(a,b)$，$(a,\bar b)$，$(\bar a,b)$，$(a,b)$。其中 $a$，$\bar a$ 分别表示甲组研发成功和失败；$b$，$\bar b$ 分别表示乙组研发成功和失败。
（1） 若某组成功研发一种新产品，则给该组记 $1$ 分，否则记 $0$ 分，试计算甲、乙两组研发新产品的成绩的平均数和方差，并比较甲、乙两组的研发水平；
（2） 若该企业安排甲、乙两组各自研发一种新产品，试估计恰有一组研发成功的概率。

\Needspace{4\baselineskip}
\textbf{解析}
\textbf{（1）} 抽取的 $15$ 次结果中，甲组成功（$a$）$10$ 次、失败（$\bar a$）$5$ 次；乙组成功（$b$）$9$ 次、失败（$\bar b$）$6$ 次。
\Needspace{5\baselineskip}
甲组成绩 $X$ 取 $1$ 共 $10$ 次、取 $0$ 共 $5$ 次，平均数与方差分别为
\nopagebreak[4]
\[\bar X=\frac{10\times1+5\times0}{15}=\frac23,\]
\nopagebreak[4]
\[s_X^2=\frac1{15}\left[10\left(1-\frac23\right)^2+5\left(0-\frac23\right)^2\right]
=\frac1{15}\left(\frac{10}9+\frac{20}9\right)=\frac29 .\]
\Needspace{5\baselineskip}
乙组成绩 $Y$ 取 $1$ 共 $9$ 次、取 $0$ 共 $6$ 次，平均数与方差分别为
\nopagebreak[4]
\[\bar Y=\frac{9\times1+6\times0}{15}=\frac35,\]
\nopagebreak[4]
\[s_Y^2=\frac1{15}\left[9\left(1-\frac35\right)^2+6\left(0-\frac35\right)^2\right]
=\frac1{15}\left(\frac{36}{25}+\frac{54}{25}\right)=\frac6{25}.\]
\Needspace{5\baselineskip}
比较：
\nopagebreak[4]
\[\bar X=\frac23>\frac35=\bar Y,\qquad s_X^2=\frac29<\frac6{25}.\]
甲组平均成绩较高，且方差较小、发挥更稳定，故甲组的研发水平较高。
\Needspace{5\baselineskip}
\textbf{（2）} 用样本频率估计概率：
\nopagebreak[4]
\[P(\text{\text{甲成功}})=\frac{10}{15}=\frac23,\qquad P(\text{\text{乙成功}})=\frac9{15}=\frac35 .\]
\Needspace{5\baselineskip}
两组各自研发，结果相互独立。“恰有一组成功”包括“甲成功且乙失败”与“甲失败且乙成功”，故
\nopagebreak[4]
\[\begin{aligned}
P&=P(\text{\text{甲成功}})P(\text{\text{乙失败}})+P(\text{\text{甲失败}})P(\text{\text{乙成功}})\\
&=\frac23\left(1-\frac35\right)+\left(1-\frac23\right)\frac35\\
&=\frac23\cdot\frac25+\frac13\cdot\frac35
=\frac4{15}+\frac3{15}=\frac7{15}.
\end{aligned}\]

\gkanswer{答案（1）：甲组平均数 $\dfrac23$、方差 $\dfrac29$；乙组平均数 $\dfrac35$、方差 $\dfrac6{25}$；甲组平均数较高且方差较小，研发水平较高。$15$ 次结果中甲成功 $10$ 次、乙成功 $9$ 次。}

\gkanswer{答案（2）：$\dfrac7{15}$。以频率估计概率，$P=\dfrac23\cdot\dfrac25+\dfrac13\cdot\dfrac35=\dfrac7{15}$。}

\end{answer}
\begin{answer}
\gksolhead{18}{}

\textbf{题目}：

如图，已知二面角 $\alpha-MN-\beta$ 的大小为 $60^\circ$，菱形 $ABCD$ 在面 $\beta$ 内，$A$，$B$ 两点在棱 $MN$ 上，$\angle BAD=60^\circ$，$E$ 是 $AB$ 的中点，$DO\perp$ 面 $\alpha$，垂足为 $O$。
\nopagebreak[4]
\begin{center}% ============================================================
% f2014_hunan_arts_q18_dihedral.tex
% 2014 湖南卷（文科）第 18 题：二面角 α-MN-β 与菱形 ABCD
% 来源 PDF：普通高考/2014/2014湖南文.pdf
% 原页：第 1 页（A3 横排）第 3 栏下半，印面 887，第 18 题题干下方
% 结构：半平面 α 在棱 MN 下方，半平面 β 在棱 MN 上方，二者成 60°；
%       菱形 ABCD 在 β 内，A、B 在 MN 上，E 为 AB 中点；
%       DO⊥α 于 O。实线：菱形四边、DE、DO；虚线：MN 上被平面 DOE 遮住的一段。
% ============================================================
\begin{tikzpicture}[>=Stealth, line width=0.7pt, font=\small,
  x=1cm, y=1cm, line cap=round, line join=round]
  % 半平面 β
  \draw (-3.60,0) -- (-4.05,2.90) -- (3.15,2.90) -- (3.60,0);
  % 半平面 α
  \draw (-3.60,0) -- (-5.20,-2.40) -- (2.00,-2.40) -- (3.60,0);
  % 棱 MN：M→A→(D 的垂足) 实线，(垂足)→E 虚线，E→B→N 实线
  \draw (-3.60,0) -- (-1.50,0);
  \draw[dashed] (-1.50,0) -- (-0.30,0);
  \draw (-0.30,0) -- (3.60,0);
  % 菱形 ABCD 的四边与对角线 DE
  \draw (-2.30,0) -- (-1.50,2.50);
  \draw (-1.50,2.50) -- (2.50,2.50);
  \draw (2.50,2.50) -- (1.70,0);
  \draw (-0.30,0) -- (-1.50,2.50);
  % DO⊥α
  \draw (-1.50,2.50) -- (-1.50,-1.05);
  % 顶点与标注
  \node[left]       at (-3.75,0.05) {$M$};
  \node[right]      at (3.75,0.05) {$N$};
  \node[below]      at (-2.30,-0.22) {$A$};
  \node[below right] at (-0.22,-0.10) {$E$};
  \node[below]      at (1.70,-0.22) {$B$};
  \node[above left] at (-1.50,2.50) {$D$};
  \node[above right] at (2.42,2.50) {$C$};
  \node[below]      at (-1.50,-1.35) {$O$};
  \node at (-3.05,2.30) {$\beta$};
  \node at (-3.95,-1.55) {$\alpha$};
\end{tikzpicture}
\end{center}
\nopagebreak[4]
（1） 证明：$AB\perp$ 平面 $ODE$；
（2） 求异面直线 $BC$ 与 $OD$ 所成角的余弦值。

\Needspace{4\baselineskip}
\textbf{解析}
\Needspace{5\baselineskip}
\textbf{（1）} 因为 $DO\perp$ 面 $\alpha$，而 $AB\subset MN\subset\alpha$，所以
\nopagebreak[4]
\[DO\perp AB .\]
\Needspace{5\baselineskip}
菱形 $ABCD$ 中 $AB=AD$，又 $\angle BAD=60^\circ$，故 $\triangle ABD$ 为正三角形。$E$ 是 $AB$ 的中点，所以
\nopagebreak[4]
\[DE\perp AB .\]
\Needspace{5\baselineskip}
$DE$ 与 $DO$ 都在平面 $ODE$ 内，且相交于 $D$，由线面垂直的判定定理得
\nopagebreak[4]
\[AB\perp\text{\text{平面} }ODE .\]
\textbf{（2）} 由（1）$AB\perp$ 平面 $ODE$，故 $AB\perp DE$，$AB\perp DO$。不妨设 $AB=2$。
\Needspace{5\baselineskip}
以 $E$ 为原点，$\overrightarrow{EB}$ 的方向为 $x$ 轴正方向；在平面 $ODE$ 内取与 $AB$ 垂直的方向为 $y$ 轴；$DO$ 的方向为 $z$ 轴，建立空间直角坐标系。则
\nopagebreak[4]
\[A(-1,0,0),\quad B(1,0,0),\quad E(0,0,0).\]
\Needspace{5\baselineskip}
因为 $AB\perp$ 平面 $ODE$，所以 $D$，$O$ 的横坐标都是 $0$。设
\nopagebreak[4]
\[D(0,y,z)\quad(z>0),\qquad O(0,y,0).\]
\Needspace{5\baselineskip}
由 $|AD|=|AB|=2$ 得
\nopagebreak[4]
\[1+y^2+z^2=4. \tag{1}\]
\Needspace{5\baselineskip}
因为 $DO\perp\alpha$，所以 $\alpha$ 是平面 $z=0$。面 $\beta$ 由 $x$ 轴与点 $D$ 确定，其法向量可取
\nopagebreak[4]
\[\boldsymbol n=(1,0,0)\times(0,y,z)=(0,-z,y).\]
\Needspace{5\baselineskip}
二面角 $\alpha-MN-\beta$ 的大小为 $60^\circ$，而 $\alpha$ 的法向量为 $(0,0,1)$，故
\nopagebreak[4]
\[\frac{|\boldsymbol n\cdot(0,0,1)|}{|\boldsymbol n|}
=\frac{|y|}{\sqrt{y^2+z^2}}=\cos60^\circ=\frac12 .\]
\Needspace{5\baselineskip}
两边平方得
\nopagebreak[4]
\[\frac{y^2}{y^2+z^2}=\frac14\Longrightarrow z^2=3y^2 .\]
\Needspace{5\baselineskip}
代入（1）：
\nopagebreak[4]
\[1+4y^2=4\Longrightarrow y^2=\frac34,\qquad z^2=\frac94,\qquad z=\frac32 .\]
\Needspace{5\baselineskip}
于是可取
\nopagebreak[4]
\[D\left(0,\frac{\sqrt3}2,\frac32\right),\qquad O\left(0,\frac{\sqrt3}2,0\right).\]
\Needspace{5\baselineskip}
从而
\nopagebreak[4]
\[\overrightarrow{AD}=\left(1,\frac{\sqrt3}2,\frac32\right),\qquad
\overrightarrow{OD}=\left(0,0,\frac32\right),\]
\nopagebreak[4]
\[\cos\langle\overrightarrow{AD},\overrightarrow{OD}\rangle
=\frac{\overrightarrow{AD}\cdot\overrightarrow{OD}}{|\overrightarrow{AD}||\overrightarrow{OD}|}
=\frac{\frac94}{2\cdot\frac32}=\frac34 .\]
\Needspace{5\baselineskip}
菱形中 $BC\parallel AD$，所以异面直线 $BC$ 与 $OD$ 所成的角等于 $\overrightarrow{AD}$ 与 $\overrightarrow{OD}$ 的夹角（取锐角），其余弦值为
\nopagebreak[4]
\[\frac34 .\]

\gkanswer{答案（1）：见解析版}

\gkanswer{答案（2）：$\dfrac34$。$BC\parallel AD$，建系求得 $\cos\langle\overrightarrow{AD},\overrightarrow{OD}\rangle=\dfrac34$。}

\end{answer}
\begin{answer}
\gksolhead{19}{}

\textbf{题目}：

如图，在平面四边形 $ABCD$ 中，$DA\perp AB$，$DE=1$，$EC=\sqrt7$，$EA=2$，$\angle ADC=\dfrac{2\pi}3$，$\angle BEC=\dfrac\pi3$。
\nopagebreak[4]
\begin{center}% ============================================================
% f2014_hunan_arts_q19_quad.tex
% 2014 湖南卷（文科）第 19 题：平面四边形 ABCD
% 来源 PDF：普通高考/2014/2014湖南文.pdf
% 原页：第 2 页（A3 横排）左栏，印面 888，第 19 题题干下方
% 结构：DA⊥AB，E 在 DA 上，DE=1，EA=2；按题设精确作图（比例 0.55 cm/单位）：
%       A(0,0)，B(6√3,0)，E(0,2)，D(0,3)，C(√3,4)。
% ============================================================
\begin{tikzpicture}[>=Stealth, line width=0.7pt, font=\small,
  x=1cm, y=1cm, line cap=round, line join=round]
  \coordinate (A) at (0,0);
  \coordinate (B) at (5.71,0);
  \coordinate (E) at (0,1.10);
  \coordinate (D) at (0,1.65);
  \coordinate (C) at (0.95,2.20);
  % 四边形边界与对角线
  \draw (A) -- (B);
  \draw (B) -- (C);
  \draw (C) -- (D);
  \draw (D) -- (A);
  \draw (E) -- (C);
  \draw (E) -- (B);
  % 顶点标注
  \node[below left]  at (A) {$A$};
  \node[below right] at (B) {$B$};
  \node[left]        at (E) {$E$};
  \node[left]        at (D) {$D$};
  \node[above]       at (C) {$C$};
\end{tikzpicture}
\end{center}
\nopagebreak[4]
（1） 求 $\sin\angle CED$ 的值；
（2） 求 $BE$ 的长。

\Needspace{4\baselineskip}
\textbf{解析}
\Needspace{5\baselineskip}
\textbf{（1）} 由 $DA\perp AB$ 及图中位置，$A$，$E$，$D$ 共线且 $E$ 在 $DA$ 上，故
\nopagebreak[4]
\[\angle CDE=\angle CDA=\frac{2\pi}3 .\]
\Needspace{5\baselineskip}
在 $\triangle CDE$ 中，$DE=1$，$CE=\sqrt7$，由余弦定理
\nopagebreak[4]
\[CE^2=DE^2+DC^2-2\cdot DE\cdot DC\cos\angle CDE,\]
\Needspace{5\baselineskip}
即
\nopagebreak[4]
\[7=1+DC^2-2\cdot1\cdot DC\cdot\left(-\frac12\right)=1+DC^2+DC .\]
\Needspace{5\baselineskip}
整理得
\nopagebreak[4]
\[DC^2+DC-6=0\Longrightarrow(DC+3)(DC-2)=0,\]
故 $DC=2$。
\Needspace{5\baselineskip}
再由正弦定理
\nopagebreak[4]
\[\frac{DC}{\sin\angle CED}=\frac{CE}{\sin\angle CDE},\]
\Needspace{5\baselineskip}
得
\nopagebreak[4]
\[\sin\angle CED=\frac{DC\sin\angle CDE}{CE}
=\frac{2\cdot\sin\frac{2\pi}3}{\sqrt7}
=\frac{2\cdot\frac{\sqrt3}2}{\sqrt7}
=\frac{\sqrt3}{\sqrt7}=\frac{\sqrt{21}}7 .\]
\Needspace{5\baselineskip}
\textbf{（2）} 由（1），$\angle CED$ 为锐角，且
\nopagebreak[4]
\[\cos\angle CED=\sqrt{1-\sin^2\angle CED}
=\sqrt{1-\frac37}=\frac2{\sqrt7}.\]
\Needspace{5\baselineskip}
因为 $A$，$E$，$D$ 共线，射线 $EA$ 与射线 $ED$ 方向相反，所以
\nopagebreak[4]
\[\angle AEC=\pi-\angle CED,\]
\nopagebreak[4]
\[\cos\angle AEC=-\cos\angle CED=-\frac2{\sqrt7},\qquad
\sin\angle AEC=\sin\angle CED=\frac{\sqrt3}{\sqrt7}.\]
\Needspace{5\baselineskip}
由图，射线 $EB$ 在 $\angle AEC$ 内部，且 $\angle BEC=\dfrac\pi3$，故
\nopagebreak[4]
\[\angle AEB=\angle AEC-\frac\pi3 .\]
\Needspace{5\baselineskip}
于是
\nopagebreak[4]
\[\begin{aligned}
\cos\angle AEB
&=\cos\angle AEC\cos\frac\pi3+\sin\angle AEC\sin\frac\pi3\\
&=\left(-\frac2{\sqrt7}\right)\cdot\frac12+\frac{\sqrt3}{\sqrt7}\cdot\frac{\sqrt3}2\\
&=-\frac1{\sqrt7}+\frac3{2\sqrt7}=\frac1{2\sqrt7}.
\end{aligned}\]
\Needspace{5\baselineskip}
在 Rt$\triangle ABE$ 中，$\angle EAB=90^\circ$，$EA=2$，故
\nopagebreak[4]
\[\cos\angle AEB=\frac{EA}{BE}\Longrightarrow
BE=\frac{EA}{\cos\angle AEB}=\frac{2}{\frac1{2\sqrt7}}=4\sqrt7 .\]

\gkanswer{答案（1）：$\dfrac{\sqrt{21}}7$。$\angle CDE=120^\circ$，由余弦定理得 $DC=2$，再由正弦定理得 $\sin\angle CED=\dfrac{2\sin120^\circ}{\sqrt7}=\dfrac{\sqrt{21}}7$。}

\gkanswer{答案（2）：$4\sqrt7$。$\cos\angle AEB=\cos\left(\angle AEC-60^\circ\right)=\dfrac1{2\sqrt7}$，在 Rt$\triangle ABE$ 中 $BE=\dfrac{EA}{\cos\angle AEB}=4\sqrt7$。}

\end{answer}
\begin{answer}
\gksolhead{20}{}

\textbf{题目}：

如图，$O$ 为坐标原点，双曲线 $C_1:\dfrac{x^2}{a_1^2}-\dfrac{y^2}{b_1^2}=1$（$a_1>0$，$b_1>0$）和椭圆 $C_2:\dfrac{y^2}{a_2^2}+\dfrac{x^2}{b_2^2}=1$（$a_2>b_2>0$）均过点 $P\left(\dfrac{2\sqrt3}3,1\right)$，且以 $C_1$ 的两个顶点和 $C_2$ 的两个焦点为顶点的四边形是面积为 $2$ 的正方形。
\nopagebreak[4]
\begin{center}% ============================================================
% f2014_hunan_arts_q20_conics.tex
% 2014 湖南卷（文科）第 20 题：双曲线 C1 与椭圆 C2
% 来源 PDF：普通高考/2014/2014湖南文.pdf
% 原页：第 2 页（A3 横排）中栏，印面 888，第 20 题题干下方
% 结构：C1: x^2-y^2/3=1（左右两支）；C2: y^2/3+x^2/2=1（竖长椭圆）；
%       以 C1 两顶点 (±1,0) 与 C2 两焦点 (0,±1) 为顶点的正方形（菱形摆放）；
%       P(2√3/3,1) 为两曲线在第一象限的交点。
% ============================================================
\begin{tikzpicture}[>=Stealth, line width=0.6pt, font=\small,
  x=1.35cm, y=1.35cm, line cap=round, line join=round]
  % 坐标轴
  \draw[->] (-2.55,0) -- (2.75,0) node[right] {$x$};
  \draw[->] (0,-2.45) -- (0,2.65) node[above] {$y$};
  % 椭圆 C2
  \draw[domain=0:360, samples=240, smooth, variable=\t]
    plot ({1.41421*cos(\t)}, {1.73205*sin(\t)});
  % 双曲线 C1 的左右两支
  \draw[domain=-2.55:2.55, samples=120, smooth, variable=\y]
    plot ({sqrt(1+\y*\y/3)}, {\y});
  \draw[domain=-2.55:2.55, samples=120, smooth, variable=\y]
    plot ({-sqrt(1+\y*\y/3)}, {\y});
  % 正方形（顶点为 (±1,0)、(0,±1)）
  \draw (1,0) -- (0,1) -- (-1,0) -- (0,-1) -- cycle;
  % 原点与点 P
  \node[below left] at (-0.06,-0.04) {$O$};
  \fill (1.15470,1) circle (1.2pt);
  \node[right] at (1.24,1.10) {$P$};
\end{tikzpicture}
\end{center}
\nopagebreak[4]
（1） 求 $C_1$，$C_2$ 的方程；
（2） 是否存在直线 $l$，使得 $l$ 与 $C_1$ 交于 $A$，$B$ 两点，与 $C_2$ 只有一个公共点，且 $|\overrightarrow{OA}+\overrightarrow{OB}|=|\overrightarrow{AB}|$？证明你的结论。

\Needspace{4\baselineskip}
\textbf{解析}
\Needspace{5\baselineskip}
\textbf{（1）} 设椭圆 $C_2$ 的半焦距为 $c_2>0$，则
\nopagebreak[4]
\[c_2^2=a_2^2-b_2^2,\]
两焦点为 $(0,\pm c_2)$。双曲线 $C_1$ 的两顶点为 $(\pm a_1,0)$。
\Needspace{5\baselineskip}
以 $(\pm a_1,0)$、$(0,\pm c_2)$ 为顶点的四边形，两条对角线长分别为 $2a_1$、$2c_2$，且互相垂直平分，故为菱形。它为正方形当且仅当两条对角线相等：
\nopagebreak[4]
\[2a_1=2c_2\Longrightarrow a_1=c_2 .\]
\Needspace{5\baselineskip}
正方形面积为
\nopagebreak[4]
\[S=\frac12\cdot2a_1\cdot2c_2=2a_1c_2=2a_1^2=2,\]
\Needspace{5\baselineskip}
故
\nopagebreak[4]
\[a_1=c_2=1 .\]
\Needspace{5\baselineskip}
点 $P\left(\dfrac{2\sqrt3}3,1\right)$ 在 $C_1$ 上，代入
\nopagebreak[4]
\[\frac{x^2}{a_1^2}-\frac{y^2}{b_1^2}=1\]
\Needspace{5\baselineskip}
得
\nopagebreak[4]
\[\frac43-\frac1{b_1^2}=1\Longrightarrow\frac1{b_1^2}=\frac13\Longrightarrow b_1^2=3 .\]
\Needspace{5\baselineskip}
故
\nopagebreak[4]
\[C_1:\ x^2-\frac{y^2}3=1 .\]
\Needspace{5\baselineskip}
点 $P$ 也在 $C_2$ 上，且 $a_2^2-b_2^2=c_2^2=1$，即 $a_2^2=b_2^2+1$。代入
\nopagebreak[4]
\[\frac{y^2}{a_2^2}+\frac{x^2}{b_2^2}=1\]
\Needspace{5\baselineskip}
得
\nopagebreak[4]
\[\frac1{b_2^2+1}+\frac{4/3}{b_2^2}=1 .\]
\Needspace{5\baselineskip}
令 $u=b_2^2>0$，去分母：
\nopagebreak[4]
\[3u+4(u+1)=3u(u+1)\Longrightarrow3u^2-4u-4=0 .\]
\Needspace{5\baselineskip}
解得
\nopagebreak[4]
\[u=\frac{4\pm\sqrt{16+48}}6=\frac{4\pm8}6,\]
\Needspace{5\baselineskip}
即 $u=2$ 或 $u=-\dfrac23$（舍）。故
\nopagebreak[4]
\[b_2^2=2,\qquad a_2^2=b_2^2+1=3,\]
\nopagebreak[4]
\[C_2:\ \frac{y^2}3+\frac{x^2}2=1 .\]
\textbf{（2）} 假设存在满足条件的直线 $l$。
\Needspace{5\baselineskip}
先化简题设等式。注意 $\overrightarrow{AB}=\overrightarrow{OB}-\overrightarrow{OA}$，由
\nopagebreak[4]
\[|\overrightarrow{OA}+\overrightarrow{OB}|=|\overrightarrow{OB}-\overrightarrow{OA}|\]
\Needspace{5\baselineskip}
两边平方得
\nopagebreak[4]
\[|\overrightarrow{OA}|^2+|\overrightarrow{OB}|^2+2\overrightarrow{OA}\cdot\overrightarrow{OB}
=|\overrightarrow{OA}|^2+|\overrightarrow{OB}|^2-2\overrightarrow{OA}\cdot\overrightarrow{OB},\]
\Needspace{5\baselineskip}
故
\nopagebreak[4]
\[\overrightarrow{OA}\cdot\overrightarrow{OB}=0,\]
即 $OA\perp OB$。
\Needspace{5\baselineskip}
若 $l$ 垂直于 $x$ 轴，设 $l:x=t$。$l$ 与 $C_2$ 只有一个公共点，需 $l$ 是椭圆的竖直切线，即 $t^2=2$。此时 $l$ 与 $C_1$ 的交点满足
\nopagebreak[4]
\[t^2-\frac{y^2}3=1\Longrightarrow y^2=3(t^2-1)=3,\]
\Needspace{5\baselineskip}
得 $A(t,\sqrt3)$，$B(t,-\sqrt3)$，于是
\nopagebreak[4]
\[\overrightarrow{OA}\cdot\overrightarrow{OB}=t^2-3=2-3=-1\neq0,\]
不合。
\Needspace{5\baselineskip}
若 $l$ 不垂直于 $x$ 轴，设
\nopagebreak[4]
\[l:\ y=kx+m .\]
\Needspace{5\baselineskip}
$l$ 与 $C_2$ 只有一个公共点。代入 $\dfrac{y^2}3+\dfrac{x^2}2=1$：
\nopagebreak[4]
\[\frac{(kx+m)^2}3+\frac{x^2}2=1
\Longrightarrow(2k^2+3)x^2+4kmx+2m^2-6=0 .\]
\Needspace{5\baselineskip}
判别式为零：
\nopagebreak[4]
\[\Delta=(4km)^2-4(2k^2+3)(2m^2-6)=0 .\]
\Needspace{5\baselineskip}
化简：
\nopagebreak[4]
\[\begin{aligned}
\Delta&=16k^2m^2-4\left(4k^2m^2-12k^2+6m^2-18\right)\\
&=48k^2-24m^2+72=24\left(2k^2-m^2+3\right)=0,
\end{aligned}\]
\Needspace{5\baselineskip}
故
\nopagebreak[4]
\[m^2=2k^2+3. \tag{1}\]
\Needspace{5\baselineskip}
$l$ 与 $C_1$ 交于两点。代入 $x^2-\dfrac{y^2}3=1$：
\nopagebreak[4]
\[3x^2-(kx+m)^2=3\Longrightarrow(3-k^2)x^2-2kmx-(m^2+3)=0 .\]
\Needspace{5\baselineskip}
若 $k^2=3$，方程变为一次方程，$l$ 与 $C_1$ 只有一个交点，不合题意，故 $k^2\neq3$。设两交点横坐标为 $x_1$，$x_2$，则
\nopagebreak[4]
\[x_1+x_2=\frac{2km}{3-k^2},\qquad x_1x_2=-\frac{m^2+3}{3-k^2}.\]
\Needspace{5\baselineskip}
由 $y_i=kx_i+m$，
\nopagebreak[4]
\[\begin{aligned}
\overrightarrow{OA}\cdot\overrightarrow{OB}
&=x_1x_2+y_1y_2\\
&=(1+k^2)x_1x_2+km(x_1+x_2)+m^2\\
&=\frac{-(1+k^2)(m^2+3)+2k^2m^2+m^2(3-k^2)}{3-k^2}\\
&=\frac{2m^2-3(1+k^2)}{3-k^2}.
\end{aligned}\]
\Needspace{5\baselineskip}
由 $\overrightarrow{OA}\cdot\overrightarrow{OB}=0$ 得
\nopagebreak[4]
\[2m^2=3(1+k^2). \tag{2}\]
\Needspace{5\baselineskip}
把（1）代入（2）：
\nopagebreak[4]
\[2(2k^2+3)=3(1+k^2)\Longrightarrow4k^2+6=3+3k^2\Longrightarrow k^2=-3,\]
无实数解。
综上，不存在满足条件的直线 $l$。

\gkanswer{答案（1）：$C_1:x^2-\dfrac{y^2}3=1$；$C_2:\dfrac{y^2}3+\dfrac{x^2}2=1$。正方形对角线相等且面积为 $2$ 给出 $a_1=c_2=1$，再代入 $P$ 定出 $b_1^2=3$，$b_2^2=2$。}

\gkanswer{答案（2）：不存在。条件等价于 $\overrightarrow{OA}\perp\overrightarrow{OB}$；与 $C_2$ 相切给出 $m^2=2k^2+3$，而 $\overrightarrow{OA}\cdot\overrightarrow{OB}=0$ 要求 $2m^2=3(1+k^2)$，联立得 $k^2=-3$，无解。}

\end{answer}
\begin{answer}
\gksolhead{21}{}

\textbf{题目}：

已知函数 $f(x)=x\cos x-\sin x+1$（$x>0$）。
（1） 求 $f(x)$ 的单调区间；
（2） 记 $x_i$ 为 $f(x)$ 的从小到大的第 $i$（$i\in\mathbf N^*$）个零点，证明：对一切 $n\in\mathbf N^*$，有 $\dfrac1{x_1^2}+\dfrac1{x_2^2}+\cdots+\dfrac1{x_n^2}<\dfrac23$。

\Needspace{4\baselineskip}
\textbf{解析}
\Needspace{5\baselineskip}
\textbf{（1）} 求导：
\nopagebreak[4]
\[f'(x)=\cos x-x\sin x-\cos x=-x\sin x .\]
因为 $x>0$，所以 $f'(x)$ 的符号由 $-\sin x$ 决定。
当 $x\in(2k\pi,(2k+1)\pi)$（$k\in\mathbf N$）时 $\sin x>0$，$f'(x)<0$，$f(x)$ 单调递减；
当 $x\in((2k+1)\pi,(2k+2)\pi)$（$k\in\mathbf N$）时 $\sin x<0$，$f'(x)>0$，$f(x)$ 单调递增。
\Needspace{5\baselineskip}
故 $f(x)$ 的单调递减区间为
\nopagebreak[4]
\[\left(2k\pi,(2k+1)\pi\right)\quad(k\in\mathbf N),\]
\Needspace{5\baselineskip}
单调递增区间为
\nopagebreak[4]
\[\left((2k+1)\pi,(2k+2)\pi\right)\quad(k\in\mathbf N).\]
\Needspace{5\baselineskip}
\textbf{（2）} 先确定零点的位置。对 $k\in\mathbf N$，
\nopagebreak[4]
\[f(k\pi)=k\pi\cos k\pi-\sin k\pi+1=k\pi(-1)^k+1 .\]
\Needspace{5\baselineskip}
于是
\nopagebreak[4]
\[f(2m\pi)=2m\pi+1>0,\qquad f((2m+1)\pi)=1-(2m+1)\pi<0 .\]
\Needspace{5\baselineskip}
对 $n\in\mathbf N^*$，区间 $((n-1)\pi,n\pi)$ 的两端点函数值异号；又 $f'(x)=-x\sin x$ 在该区间内不变号，$f$ 严格单调，故其中恰有一个零点。因此
\nopagebreak[4]
\[(n-1)\pi<x_n<n\pi .\]
\Needspace{5\baselineskip}
特别地，$x_1\in(0,\pi)$。由
\nopagebreak[4]
\[f\left(\frac\pi2\right)=\frac\pi2\cos\frac\pi2-\sin\frac\pi2+1=0-1+1=0\]
\Needspace{5\baselineskip}
及 $f$ 在 $(0,\pi)$ 上严格递减，得
\nopagebreak[4]
\[x_1=\frac\pi2 .\]
\Needspace{5\baselineskip}
对 $i\geqslant2$，由 $(i-1)\pi<x_i$ 得
\nopagebreak[4]
\[\frac1{x_i^2}<\frac1{(i-1)^2\pi^2} .\]
\Needspace{5\baselineskip}
于是
\nopagebreak[4]
\[\sum_{i=1}^n\frac1{x_i^2}
<\frac4{\pi^2}+\frac1{\pi^2}\sum_{j=1}^{n-1}\frac1{j^2} .\]
\Needspace{5\baselineskip}
当 $j\geqslant2$ 时，
\nopagebreak[4]
\[\frac1{j^2}<\frac1{j(j-1)}=\frac1{j-1}-\frac1j,\]
\Needspace{5\baselineskip}
故
\nopagebreak[4]
\[\sum_{j=1}^{n-1}\frac1{j^2}
<1+\sum_{j=2}^{n-1}\left(\frac1{j-1}-\frac1j\right)
=1+\left(1-\frac1{n-1}\right)<2 .\]
\Needspace{5\baselineskip}
所以
\nopagebreak[4]
\[\sum_{i=1}^n\frac1{x_i^2}<\frac4{\pi^2}+\frac2{\pi^2}=\frac6{\pi^2} .\]
\Needspace{5\baselineskip}
由 $\pi>3$ 得 $\pi^2>9$，从而
\nopagebreak[4]
\[\frac6{\pi^2}<\frac69=\frac23 .\]
\Needspace{5\baselineskip}
故对一切 $n\in\mathbf N^*$，
\nopagebreak[4]
\[\frac1{x_1^2}+\frac1{x_2^2}+\cdots+\frac1{x_n^2}<\frac23 .\]

\gkanswer{答案（1）：单调递减区间为 $(2k\pi,(2k+1)\pi)$，单调递增区间为 $((2k+1)\pi,(2k+2)\pi)$，$k\in\mathbf N$。$f'(x)=-x\sin x$。}

\gkanswer{答案（2）：见解析版}

\end{answer}